\ifdefined\pdftrailerid\pdftrailerid{}\fi
\documentclass[10pt,a4paper,oneside,openany]{report}
\usepackage[T1]{fontenc}
\usepackage[utf8]{inputenc}
\DeclareUnicodeCharacter{0393}{\ensuremath{\Gamma}}
\DeclareUnicodeCharacter{2082}{\ensuremath{{}_{2}}}
\usepackage{lmodern,amsmath,amssymb,amsthm,mathtools,mathrsfs}
\usepackage{array,longtable,booktabs,calc,graphicx,enumitem,needspace}
\usepackage[margin=22mm,top=21mm,bottom=23mm]{geometry}
\usepackage{xcolor,hyperref}
\definecolor{ink}{RGB}{24,60,79}
\hypersetup{colorlinks=true,linkcolor=ink,urlcolor=ink,citecolor=ink,
  bookmarksdepth=1,pdftitle={A human-readable version of the formal Alpöge–Claude proof of RH and Dirichlet GRH — Detailed edition}}
\urlstyle{same}
\raggedbottom
\setcounter{tocdepth}{1}
\setcounter{secnumdepth}{1}
\setlength{\parindent}{0pt}
\setlength{\parskip}{4pt plus 1pt}
\setlength{\emergencystretch}{3em}
\widowpenalty=8000
\clubpenalty=8000
\newcommand{\tightlist}{\setlength{\itemsep}{0pt}\setlength{\parskip}{0pt}}
\newcommand{\passthrough}[1]{#1}
\newtheorem{theorem}{Theorem}[section]
\newtheorem{proposition}[theorem]{Proposition}
\newtheorem{lemma}[theorem]{Lemma}
\newtheorem{corollary}[theorem]{Corollary}
\theoremstyle{remark}\newtheorem{remark}[theorem]{Remark}
\newcommand{\Rea}{\operatorname{Re}}
\newcommand{\HS}{\mathrm{HS}}
\newcommand{\tr}{\operatorname{tr}}
\newcommand{\codim}{\operatorname{codim}}
\newcommand{\covol}{\operatorname{covol}}
\newcommand{\OK}{\mathcal O_K}
\providecommand{\Sha}{\operatorname{Sha}}
\newcommand{\src}[2]{\href{https://github.com/CocoSgt/riemann-hypothesis-lean/blob/7cbf7a0c32fca1badf05931879a37c3f8a2b03b4/src/#1}{\texttt{#2}}}
\newcommand{\status}[1]{\par\textit{Scope recorded in the original derivation.} #1\par}
\begin{document}
\begin{titlepage}
\vspace*{16mm}
{\color{ink}\LARGE\bfseries A human-readable version of the formal Alpöge–Claude proof of RH and Dirichlet GRH\par}
\vspace{5mm}
{\Large Detailed edition\par}
\vspace{3mm}
@EDITION@ \quad @DATE@\par
\vspace{10mm}
For every Dirichlet character, the zeros of its $L$-function in the open critical strip lie on the line $\Re s=1/2$. The character of conductor one gives the assertion for the Riemann zeta function. This volume follows the proof of Levent Alpöge and Claude from the arithmetic construction to these conclusions.
\par\medskip
The first chapter follows the argument from theta derivatives to a bound on weighted prime sums. The later chapters derive the estimates used at each stage. Citations to the formal archive locate the exact statements. We also credit the mathematical expositions we follow, including Kaiwen Jin's argument for detecting a zero from a one-sided estimate.
\par\medskip
\textbf{How much of the proof is here?} The manuscript derives the global Galois attachment and its ramified descent, and constructs the integral cyclotomic lattice used in the congruence argument. On the geometric tower it proves \hyperref[claim:GRT9]{equality of the component volumes} and constructs the \hyperref[claim:GRT11]{unlabelled boundary lattice}. The labelled flags at ramified places and the prescribed height estimates still require a local argument. For the real periods it computes the \hyperref[claim:FBE9]{singleton normalization} and the \hyperref[claim:FBE11]{leading signed coefficient on regular blocks}. Completing this passage requires identifying the original differential operators on every residual slice and extending their pairings through the singular blocks. The final chapter records these remaining hypotheses.
\vfill
\href{https://github.com/bmatschke/grh-alpoge-claude/tree/@REVISION@}{Source revision \texttt{@REVISION@}}\par
\href{https://github.com/CocoSgt/riemann-hypothesis-lean/tree/7cbf7a0c32fca1badf05931879a37c3f8a2b03b4}{Formal source, revision \texttt{7cbf7a0c32fc}}\par
\end{titlepage}
\pagenumbering{roman}
\begingroup
\small\setlength{\parskip}{0pt}\tableofcontents
\endgroup
\clearpage\pagenumbering{arabic}
